


\section{Fifteenth Sunday after Trinity: A Dwelling-Place in Heaven}


\begin{quote}

Matthew 6:19-23. Do not lay up for yourselves treasures on earth, where moth and rust corrupt, and where thieves break through and steal; but lay up for yourselves treasures in Heaven, where neither moth nor rust corrupts, and where thieves do not break through nor steal. For where your treasure is, there your heart will be also. The eye is the body’s light; therefore, if your eye is pure, your whole body will be full of light; but if your eye is evil, your whole body will be dark. If therefore the light that is in you is darkness, how great is that darkness!

\end{quote}

\bigskip

Our text speaks of two things: heavenly treasure and a pure eye; and yet these two are most closely connected. For together both sayings mean this, that a liberated man must have the goal of all his heart’s longings and strivings in Heaven, and that he must always look straight toward the goal without letting his sight be confused by any earthly thing. In other words: it is singleness of mind and singleness of sight that the Savior requires of His disciples.

It is the irresistible bent of the natural heart to love the world, or the things that are in the world. In one way or another, in one form or another, the world with its perishable things is the abiding place of the natural human heart. For some, it is money and possessions that draw the heart to themselves; for others, power; for others, honor; for others, pleasure; for others, earthly comfort and well-being; but no natural heart can avoid having the goal of its longings in earthly and perishable things. For the natural eye sees no farther than earthly treasures, and when the eye is dark and blind, then the whole man also remains in darkness and blindness.

Therefore, when the Lord exhorts us to gather heavenly treasures and direct the heart’s longings toward Heaven, it is an exhortation to a whole and complete conversion from the natural mind and its longings and strivings. He would place new light in our eye and set a new goal for our life, that we should not walk in darkness and perish with the world’s perdition.

The Savior sees how fleeting the world and its lust are; He sees how short men’s stay upon earth is, and how soon their life is at an end; He knows that the heart which clings to temporal and earthly things cannot become blessed when it loses them and is suddenly torn out into eternity and there must lack everything that was the heart’s delight and joy. Therefore He exhorts so earnestly to a thorough conversion from a worldly mind, and warns so heartily against letting the eye be dazzled by the vanity of earth. Therefore the whole Holy Scripture is full of exhortations to let go of the earthly and direct eye, mind, and heart toward the heavenly.

In a most special way Scripture warns against gathering treasures on earth and making these the goal of life. Therefore we read: "Those who will be rich fall into temptation and a snare, and into many foolish and hurtful desires, which sink men into destruction and perdition; for the love of money is a root of all evil; because some longed after it, they wandered from the faith and have pierced themselves through with many sorrows" (1 Tim. 6:9, 10).

It blinds the eye and stupefies and deceives the heart when a man sees no farther than the world and its lust. The eye grows dark and the mind earthly; and the whole man becomes a slave of Mammon, who with all his toil prepares for himself only an eternal misery; for his happiness was in the earthly and temporal, and when the world passes away with its lust, then there remains only an eternal deprivation, an infinite emptiness, an unsatisfied longing; and the heart suffers for ever without relief, without consolation.

Therefore, if you would become blessed, your heart must first be loosed from the world and your eye must look higher up and farther on than the vanity of the world. You must learn to love, seek, and desire the heavenly and eternal things, if your heart is to find blessedness in God’s Heaven.

How can this come to pass? It happens only by a thorough conversion. That is the only way on which a blind man of the world can receive sight; it happens only by the miracle that opens eye and heart to sin and grace, to the vanity of the world and the blessedness of Heaven. It happens only by a heartfelt repentance that my heart has fastened itself to the creature instead of the Creator, and by a humble surrender to Jesus’ cross and blood. Only the forgiveness of our sins and God’s mercy toward our poor, wandering hearts can again open our eyes to Heaven’s glory and place in our souls a longing for our true home in the Father’s house.

Have your eyes been opened in this way? Has this incomprehensible miracle taken place in you? You, who before saw only the world and its lust, have you also seen the world’s perdition? Have you cast an eye down into the abyss that opens its jaws to swallow those who so carelessly walk on along the world’s way and run so eagerly after its seductive joys? Have you, terrified and broken over your sin, taken refuge at Golgotha and from the cross looked up to the Father’s heart, which opens itself to all who come to His Son? Blessed are you then; for you have been delivered from the shameful chains of bondage with which vanity had bound your immortal soul. Blessed are you then; for you have found the Father’s heart and the Savior’s love, and you have been given to taste the goodness of God and the powers of the world to come.

But he who through Christ’s cross has died to the world and is alive to God is loosed from earth and joined to Heaven; his eye has become pure and his heart heavenly. He sees rightly, and the longings of his heart are heavenly. For him it is a daily question how he may come safely through the perilous world, how he may become more and more familiar with Heaven and its joys. It is no longer on earth, but in Heaven, that he seeks his treasure and lays up store for himself.

Is it then truly possible to gather treasures in Heaven? Is there any way by which we here in time may live for eternity, labor for eternity, lay up a treasure in eternity? Yes, by God’s wondrous grace. Those who have the opened eye and the opened heart for the incorruptible and heavenly things, those who long thither, where the Father’s home is opened for them, they also have a way and a means by which they can gather treasures in Heaven.

You live in the midst of immortal men. You dwell upon an earth that shall perish; but the human hearts that beat so near your own, they shall continue to live when the earth is no more. And if it is a glory on earth to possess the love of a heart, then it is not less, but greater, glory to possess the love of a heart in Heaven’s incorruptibility. Therefore the art of gathering treasures in Heaven consists in this, that you bind the love of human hearts by heartfelt goodness. Hear what Scripture says concerning the rich: "Charge those who are rich in this present world that they be not high-minded, nor trust in uncertain riches, but in the living God, who gives us richly all things to enjoy; that they do good, become rich in good works, gladly give, distribute, thus gathering for themselves a treasure, a good foundation for what is to come, that they may lay hold on eternal life."

But now the poor man in this world? Is it not granted to him to gather treasures in Heaven? It would indeed be hard if he for whom earthly life is so heavy and straitened should not also be able to have this joy, to gather treasures in Heaven. God be praised, there is a way just as simple for the poor as for the rich. If we gather treasures in Heaven by our love toward men, then there is room for the poor no less than for the rich. The rich need the poor man’s heart no less than the poor need the rich man’s heart. The rich can do good with their possessions; the poor can do good with their goodness, their patience, their kindness, their love. The rich are so surrounded by envy and hatred and ill will and bitterness and faithlessness, that it is a singular kindness toward them when a poor man loves them and thanks them, rejoices in their prosperity, and sees in it also a wise leading of God and a gracious gift. Heart needs heart, and it is just as easy for a poor man to gladden a human heart as it is for a rich man.

Turn, then, heart and eye toward Heaven; let the light of God’s love illumine you, and you shall learn what it is to have a heavenly treasure and to increase it from day to day.

